% !PRINTED 538
\weberpage{538}
%<<CONT
ordnung und von ihrem Anfangselement, so ergiebt sich, dass durch Anwendung irgend einer Permutation $\pi_1$ aus $P$ die Elemente der einzelnen Cyklen $\gamma$ nicht von einander getrennt, sondern nur unter einander (cyklisch) vertauscht und ausserdem die Cyklen mit einander vertauscht werden können. Die Gruppe ist also, wenn $s>1$ ist, \gesperrt{imprimitiv}.

Eine rationale Function der Argumente $\alpha,\alpha_1,\ldots\alpha_{r-1}$, die ihren Werth nicht ändert, wenn die $\alpha$ der cyklischen Permutation $\gamma$ und ihren Wiederholungen unterworfen werden, heisst eine \gesperrt{cyklische Function} der $\alpha$. Bezeichnet man mit
\weq{4}{\omega = \psi(\alpha,\alpha_1\ldots\alpha_{r-1})}
eine solche cyklische Function und mit
\weq{5}{\omega,\ \omega_1,\ \omega_2\ldots\omega_{s-1}}
die conjugirten Werthe von $\omega$, setzt also z.\,B.
\weq{6}{\omega_1 = \psi(\beta,\beta_1\ldots\beta_{r-1}),}
so sind diese Grössen nach \S.\,158 die Wurzeln einer irreducibeln Gleichung $s^{\text{ten}}$ Grades:
\weq{7}{\Phi(t) = 0,}
deren Gruppe man erhält, wenn man die durch $P$ hervorgerufenen Permutationen der Grössen (5) aufsucht.

Man erhält aber die durch $\pi_1\pi_2$ unter den Grössen (5) bewirkte Permutation, wenn man die beiden durch $\pi_1$ und $\pi_2$ einzeln hervorgerufenen Permutationen zusammensetzt, und die Gruppe der Permutationen von (5) ist daher auch commutativ. Daher ist $\Phi(t)=0$ eine \gesperrt{Abel'sche} Gleichung $s^{\text{ten}}$ Grades.

Adjungirt man $\omega$, so zerfällt $F(x)$ in $s$ Factoren $r^{\text{ten}}$ Grades:
\[
F(x) = F(x,\omega)F(x,\omega_1)\ldots F(x,\omega_{s-1}),
\]
von denen der erste, $F(x,\omega)$, die Wurzeln $\alpha,\alpha_1\ldots\alpha_{r-1}$ hat, und die Gruppe der Gleichung $F(x,\omega)=0$ besteht allein aus der Periode der cyklischen Permutation $\gamma$.

Wir wollen eine Gleichung, deren Gruppe aus einem einzigen Cyklus und seinen Wiederholungen besteht, eine \gesperrt{cyklische Gleichung} nennen, so dass die cyklischen Gleichungen der einfachste Specialfall der \gesperrt{Abel'schen} Gleichungen sind. Wir haben dann also bewiesen, dass die Lösung jeder \gesperrt{Abel'schen} Gleichung zurückgeführt wird auf die Lösung einer \gesperrt{Abel'schen} Gleichung niedrigeren Grades und auf die Lösung einer Reihe
%CONT>>
